% ==================================================================
% id:       seq-030
% topic:    סדרות
% subtopic: הגדרת הגבול (ε, N)
% solution: full
% ==================================================================
\textbf{תהי
$(a_n)_{n=1}^{\infty}$
סדרה המקיימת
$a_n\to L$.
הוכיחו או הפריכו:}
\begin{enumerate}
\item
אם קיים
$M$
כך ש-$b_n=a_n$
לכל
$n>M$,
אז
$b_n\to L$.
\hint{כתבו את הגדרת הגבול עבור $(a_n)$. החל מאיזה מקום ההגדרה מתקיימת גם עבור $(b_n)$?}
\begin{solbox}
הטענה נכונה. יהי
$\varepsilon>0$.
מכיוון ש-$a_n\to L$, קיים
$N$
כך ש-$|a_n-L|<\varepsilon$
לכל
$n>N$.
לכל
$n>\max\{N,M\}$
מתקיים
$b_n=a_n$,
ולכן
$$|b_n-L|=|a_n-L|<\varepsilon .$$
מצאנו מקום שהחל ממנו
$|b_n-L|<\varepsilon$,
ולכן
$b_n\to L$.

כלומר, שינוי של מספר סופי של איברים לא משנה את ההתכנסות ואת הגבול.
\end{solbox}

\item
אם
$b_n=a_n$
עבור אינסוף ערכים של
$n$,
אז
$b_n\to L$.
\hint{מה ההבדל בין ``לאינסוף ערכים של $n$'' לבין ``לכל $n$ החל ממקום מסוים''? נסו לחשוב על סדרה שלא מתכנסת ועל ``תיקון'' שלה.}
\begin{solbox}
הטענה לא נכונה. ניקח
$a_n=1$
לכל
$n$
(סדרה קבועה, ולכן
$a_n\to1$),
ו-$b_n=(-1)^n$.
לכל
$n$
זוגי מתקיים
$b_n=1=a_n$,
כלומר
$b_n=a_n$
לאינסוף ערכים של
$n$.
אבל הסדרה
$(-1)^n$
אינה מתכנסת.

ההבדל מהסעיף הקודם: כאן הסדרות מסכימות באינסוף מקומות, אבל לא בכל המקומות החל ממקום מסוים.
\end{solbox}

\item
$a_{n+1}\to L$.
\hint{קחו את ה-$N$ שהגדרת הגבול נותנת עבור $(a_n)$. אם $n>N$, מה אפשר לומר על $n+1$?}
\begin{solbox}
הטענה נכונה. יהי
$\varepsilon>0$,
ויהי
$N$
כך ש-$|a_n-L|<\varepsilon$
לכל
$n>N$.
אם
$n>N$,
אז גם
$n+1>N$,
ולכן
$$|a_{n+1}-L|<\varepsilon .$$
לכן
$a_{n+1}\to L$.

בהמשך תראו סוג של שאלות שבו משתמשים בנימוק הזה לא מעט.
\end{solbox}
\end{enumerate}
